% Revize 2026-10-28: obsah, matematika a návaznost; podrobnosti a časy
% viz planovani/REVIZE_PRIPRAV_1_ROCNIK.md.
% Aktuální porovnání se skutečným obsahem místních PDF Realisticky.cz:
% externi_zdroje/realisticky/01_Zakladni_poznatky/02_Ciselne_obory/010209_Usmernovani_zlomku/010209_Lekce.pdf
% Revize 2026-10-27: cílené doplnění; nové úlohy a řešení jsou vlastní tvorba.
% Metodický vzor: 4_rocnik/pripravy/Uvod_do_kombinatoriky.pdf (samostatný dokument).
% Hodina 36, 1. ročník SPŠ. Rozsah: Rozpis_vyuky_MAT_1_rocnik.tex, řádek 36.

% Plán: 1_rocnik/tematicke_plany/SPS.pdf; společná projektová šablona.

% Učebnice: ucebnice/01_Zakladni_poznatky/Mocniny_a_odmocniny/Mocniny a Odmocniny.tex

% Původ úloh: P1: přesný výraz z kostry učitele. Ostatní vlastní zadání a řešení.

% Metodický zdroj: externi_zdroje/realisticky/01_Zakladni_poznatky/02_Ciselne_obory/010209_Usmernovani_zlomku/010209_Lekce.pdf

% URL: http://www.realisticky.cz/ucebnice/01%20Matematika%20S%C5%A0/01%20Z%C3%A1kladn%C3%AD%20poznatky/02%20%C4%8C%C3%ADseln%C3%A9%20obory/09%20Usm%C4%9Br%C5%88ov%C3%A1n%C3%AD%20zlomk%C5%AF.pdf

\documentclass[11pt,a4paper]{article}

% Spolecny styl pro ucitelske pripravy MAT (Work/MAT)
% Verze 1.1, 2026-10-23
\NeedsTeXFormat{LaTeX2e}
\ProvidesPackage{priprava}[2026/10/23 v1.1 Sablona ucitelskych priprav MAT]

% Zakladni typografie a cestina
\RequirePackage[T1]{fontenc}
\RequirePackage[utf8]{inputenc}
\RequirePackage[czech]{babel}
\RequirePackage{lmodern}
\RequirePackage{microtype}
\RequirePackage{amsmath,amssymb}
\RequirePackage{xcolor}
\RequirePackage{enumitem}
\RequirePackage[most]{tcolorbox}
\RequirePackage{tabularx,booktabs,array}
\RequirePackage{fancyhdr}
\RequirePackage{needspace}
\RequirePackage{geometry}
\RequirePackage[hidelinks]{hyperref}

\geometry{a4paper,top=17mm,bottom=18mm,left=23mm,right=23mm}
\definecolor{notegray}{gray}{0.39}
\setlength{\parindent}{0pt}
\setlength{\parskip}{2.5pt plus 1pt minus 0.5pt}
\setlist[itemize]{topsep=2pt,itemsep=1pt,parsep=0pt,leftmargin=1.45em}
\setlist[enumerate]{topsep=2pt,itemsep=2pt,parsep=0pt,leftmargin=1.65em}
\setcounter{secnumdepth}{0}
\renewcommand{\arraystretch}{1.12}

% Zahlavi: vlevo rocnik | globalni tema, vpravo cislo: nazev hodiny
% Zapati zustava zcela prazdne.
\pagestyle{fancy}
\fancyhf{}
\renewcommand{\headrulewidth}{0.35pt}
\renewcommand{\footrulewidth}{0pt}
\setlength{\headheight}{14pt}
\newcommand{\nastavPripravu}[4]{%
  \fancyhead[L]{\small #1\enspace\textbar\enspace #2}%
  \fancyhead[R]{\small #3: #4}%
}

% Minimalisticke cernobile boxy. Kazda definice zvlast.
\tcbset{commonlesson/.style={enhanced,breakable=false,colback=white,colframe=black,
  colbacktitle=white,coltitle=black,boxrule=0.45pt,
  arc=0.8mm,outer arc=0.8mm,left=3mm,right=3mm,top=1.6mm,bottom=1.6mm,
  toptitle=1mm,bottomtitle=0.7mm,fonttitle=\bfseries,
  attach boxed title to top left={xshift=2.5mm,yshift=-1.15mm},
  boxed title style={colback=white,colframe=white,boxrule=0pt,left=1mm,right=1mm}}}
\newtcolorbox{zadaniBox}[1][]{commonlesson,title={Zadání},#1}
\newtcolorbox{prikladBox}[1][]{commonlesson,title={Příklad},#1}
\newtcolorbox{definiceBox}[1][]{commonlesson,title={Definice},#1}

% Spolecny box pro zadani a jeho reseni, bez mezery mezi castmi.
% Pouziti:
%   \begin{prikladReseniBox}
%     Zadani ulohy.
%     \tcblower
%     \begin{align*} ... \end{align*}
%   \end{prikladReseniBox}
% Spodni cast ma jemny sedy podklad a delici linku; lze v ni sazet
% vice radku matematiky. Obe casti maji jen jeden spolecny ramecek.
\tcbset{solutionlesson/.style={
  commonlesson,
  enhanced jigsaw,
  breakable,
  colbacklower=black!3,
  segmentation style={solid,draw=black!25,line width=0.35pt},
  fontlower=\small\color{notegray},
  before lower={\noindent\textbf{Řešení:}\enspace},
  leftlower=3mm,rightlower=3mm
}}
\newtcolorbox{prikladReseniBox}[1][]{solutionlesson,title={Příklad},#1}
\newtcolorbox{zadaniReseniBox}[1][]{solutionlesson,title={Zadání},#1}

% Pomucky pro autora pripravy: sedive poznamky; \reseni pro starsi pripravy
\newcommand{\poznamka}[1]{\par\begingroup\small\color{notegray}\itshape
  \textnormal{\textbf{Poznámka pro učitele.}}\enspace #1\par\endgroup}
\newcommand{\reseni}[1]{\par\begingroup\small\color{notegray}
  \textbf{Očekávané řešení:}\enspace #1\par\endgroup}
\newcommand{\blok}[1]{\needspace{5\baselineskip}\subsection*{#1}}
\newcommand{\vyrok}[1]{\uv{#1}}


\nastavPripravu{1. ročník}{Odmocniny}{36}{Usměrňování sdruženým výrazem}

\begin{document}

\blok{Výklad a souvislosti}

Jmenovatel $1+\sqrt5$ obsahuje součet; násobení samotnou odmocninou ho neusměrní. Vyzkoušíme změnit znaménko mezi oběma členy a při násobení sledujeme, které členy se vyruší.

\begin{definiceBox}[title={Definice},breakable=false]
\textbf{Sdružený výraz} k dvojčlenu $u+v$ je $u-v$ a naopak. Při usměrňování volíme členy tak, aby součin odstranil odmocniny ve jmenovateli.
\end{definiceBox}


Přímým roznásobením dostáváme $(u+v)(u-v)=u^2-uv+vu-v^2=u^2-v^2$. Zde tento vztah používáme jen jako prostředek usměrnění; soustavný výklad vzorců pro mnohočleny následuje v 49. hodině. Rozšiřující výraz musí být nenulový.

\blok{Řešené příklady}

\begin{prikladReseniBox}[title={Příklad 1},breakable=false]
Usměrni $\frac5{1+\sqrt5}$.
\tcblower
Jmenovatel je kladný a $1-\sqrt5\ne0$. Proto
\[\frac5{1+\sqrt5}=\frac{5(1-\sqrt5)}{(1+\sqrt5)(1-\sqrt5)}=\frac{5(1-\sqrt5)}{-4}=\frac{5(\sqrt5-1)}4.\]
\end{prikladReseniBox}


\begin{prikladReseniBox}[title={Příklad 2},breakable=false]
Usměrni $\frac1{\sqrt x-2}$ a urči podmínky.
\tcblower
Původní výraz vyžaduje $x\ge0$, $x\ne4$. Rozšiřujeme $\sqrt x+2$, které je na tomto oboru kladné. Výsledek je $\frac{\sqrt x+2}{x-4}$ při stejných podmínkách.
\end{prikladReseniBox}


\blok{Samostatné procvičení}

\begin{zadaniReseniBox}[title={Procvičení},breakable=false]
\begin{enumerate}[label=\alph*)]
\item Usměrni $\frac1{\sqrt3-1}$.
\item Usměrni $\frac2{\sqrt5+\sqrt3}$.
\item Usměrni $\frac3{2-\sqrt3}$.
\end{enumerate}
\tcblower
\begin{enumerate}[label=\alph*)]
\item $\frac{\sqrt3+1}{2}$, protože součin sdružených jmenovatelů je $3-1=2$.
\item $\sqrt5-\sqrt3$, protože po rozšíření vznikne jmenovatel $5-3=2$.
\item $3(2+\sqrt3)=6+3\sqrt3$; jmenovatel po rozšíření je $4-3=1$.
\end{enumerate}
\end{zadaniReseniBox}


\poznamka{Sleduj znaménko vzniklého jmenovatele. Rozšiřuj celý zlomek, nejen jeho část. Pro menší část třídy stačí číselné příklady; úlohu s proměnnou řeš společně až po nich.}


\blok{Navazující procvičování ve třídě}

\begin{zadaniReseniBox}[title={Rezerva: dvě odmocniny a kontrola znaménka},breakable=false]
\begin{enumerate}[label=\alph*)]
\item Usměrni $\frac3{\sqrt7+2}$ a $\frac1{\sqrt6-\sqrt2}$.
\item Vysvětli, proč $\frac1{2+\sqrt x}=\frac{2-\sqrt x}{4-x}$ nelze bez výjimky používat pro všechna $x\ge0$.
\end{enumerate}
\tcblower
\begin{enumerate}[label=\alph*)]
\item $\sqrt7-2$; $\frac{\sqrt6+\sqrt2}{4}$. Součiny sdružených výrazů jsou $3$ a $4$.
\item Pro $x=4$ je původní výraz $1/4$, ale rozšiřující výraz je nula a nový zlomek neexistuje. Vzorec platí pro $x\ge0$, $x\ne4$; hodnotu při $4$ uvedeme zvlášť.
\end{enumerate}
\end{zadaniReseniBox}

\blok{Volitelná rezerva}

\begin{prikladReseniBox}[title={Rozšíření procvičení},breakable=false]
Pro $x\ge0$ usměrni $\frac1{1+\sqrt x}$ tak, aby nezmizel přípustný případ $x=1$.
\tcblower
Pro $x\ne1$ dostaneme $\frac{1-\sqrt x}{1-x}$. Pro $x=1$ původní výraz dává $1/2$, ale rozšiřující výraz $1-\sqrt x$ by byl nulový. Proto tento případ uvedeme zvlášť; původní obor je všechna $x\ge0$.
\end{prikladReseniBox}



\poznamka{Původní výklad i příklady zůstávají základem. Nová skupina rozšiřuje zásobu rezerv, nezvyšuje povinný počet postupů; pro kontrolu lze výsledek vynásobit původním jmenovatelem. Domácí čtveřice graduje od čísel k podmínkám s proměnnou. Původní případ nulového sdruženého výrazu řeš až po upevnění běžných příkladů.}

\blok{Domácí procvičení}

\begin{zadaniReseniBox}[title={Domácí úlohy},breakable=false]
\begin{enumerate}[label=\alph*)]
\item Usměrni $\frac2{\sqrt3+1}$.
\item Usměrni $\frac4{3-\sqrt5}$.
\item Usměrni $\frac3{\sqrt7-\sqrt2}$.
\item Urči podmínky a usměrni $\frac1{\sqrt x-3}$.
\end{enumerate}
\tcblower
\begin{enumerate}[label=\alph*)]
\item $\sqrt3-1$.
\item $3+\sqrt5$.
\item $\frac{3(\sqrt7+\sqrt2)}5$.
\item $x\ge0$, $x\ne9$; výsledek $\frac{\sqrt x+3}{x-9}$. Rozšiřující výraz je vždy kladný.
\end{enumerate}
\end{zadaniReseniBox}

\end{document}
